\subsection{波兰空间中的贫集等价关系}

引理 3. --- 设 \((\mathrm{U}_{i})_{i\in\mathrm{I}}\) 是拓扑空间 X 的非空开集的一个有限族，O 是 \(X\times X\) 的一个开稠密子集。存在 X 的非空开子集的一个族 \((\mathrm{V}_{i})_{i\in\mathrm{I}}\)，使 \(V_{i}\subset U_{i}\)（对每个 \(i\in I\)），且 \(V_{i}\times V_{i'}\subset O\)（对 I 中两个不同元素的每个偶 \((i,i')\)）。

对 I 中两个不同元素的每个偶 \((j_{1}, j_{2})\)，由族 \((x_{i}) \in \mathrm{X}^{\mathrm{I}}\)（满足 \((x_{j_{1}}, x_{j_{2}}) \in \mathrm{O}\)）所成的集合是 \(X^{I}\) 的一个开稠密子集。因此，这些开集之交 \(\Omega\)——当 \((j_{1}, j_{2})\) 取遍 I 的不同元素之偶所成的有限集时——是 \(X^{I}\) 的一个开稠密集。

因此，\(\Omega \cap \prod_{i\in \mathbf{I}}\mathrm{U}_i\) 是 \(\mathrm{X}^{\mathrm{I}}\) 的一个非空开集，从而包含 \(\mathrm{X}^{\mathrm{I}}\) 的一个形如 \(\prod_{i\in \mathbf{I}}\mathrm{V}_i\) 的开集，其中对每个 \(i\in \mathbf{I}\)，\(\mathrm{V}_i\) 是 \(\mathrm{U}_i\) 的一个非空开子集。这就证明了引理。

命题 8. --- 设 P 是一个非空波兰拓扑空间，R 是 P 上的一个等价关系，其图象是 P × P 的一个贫子集。存在一个从集合 \(\{0,1\}^{\mathbf{N}}\)（赋以诸离散拓扑的乘积拓扑）到 P 中的连续单射映射，其像与 R 的每个等价类至多交于一点。

给空间 P 赋以一个与其拓扑相容且使 P 完备的度量 \(d\)。设 \((\mathrm{A}_n)_{n\in \mathbf{N}}\) 是 \(\mathrm{P}\times \mathrm{P}\) 的闭子集的一个递增序列，其内部皆为空，使得关系 R 的图象 \(\Gamma_{\mathrm{R}}\) 包含于诸 \(\mathrm{A}_n\) 之并。

设 \(\mathcal{O}\) 是 P 的非空开子集所成的集合。我们将用归纳法构造一个序列 \((f_n)_{n \in \mathbf{N}}\)，其中对每个 \(n\)，\(f_n\) 是从 \(\{0,1\}^n\) 到 \(\mathcal{O}\) 中的一个映射，它满足下列性质：

\begin{enumerate}
\def\labelenumi{(\roman{enumi})}
\item
  对每个 \(n \geqslant 1\)、每个 \(x \in \{0,1\}^{n-1}\) 及每个 \(t \in \{0,1\}\)，开集 \(f_n(x,t)\) 的闭包包含于 \(f_{n-1}(x)\)；
\item
  对每个 \(x \in \{0,1\}^n\)，有 \(\mathrm{diam}(f_n(x)) \leqslant 2^{-n}\)；
\item
  对每个 \(n \geqslant 1\) 及两个不同元素的每个偶 \((x, x')\)（属于 \(\{0, 1\}^n\)），\(f_n(x) \times f_n(x')\) 与 \(\mathrm{A}_{n-1}\) 不相交。
\end{enumerate}

选取 P 的一个直径 \(\leqslant 1\) 的非空开集 U，并把 \(f_{0}\) 定义为以 \(\{U\}\) 为像的常值映射。假设映射 \(f_{0}, \ldots, f_{n}\) 已经构造好。

设 \(p \colon \{0,1\}^{n+1} \to \{0,1\}^n\) 是由 \(p(x_0, \ldots, x_n) = (x_0, \ldots, x_{n-1})\) 定义的映射。把上述引理 3 应用于开集族 \((f_n(p(x)))\)（对 \(x \in \{0,1\}^{n+1}\)）及开稠密集 \(\mathbb{C}\mathrm{A}_n\)（\(\mathrm{P} \times \mathrm{P}\) 的开稠密集），可知存在 P 的非空开子集的一个族 \((g(x))_{x \in \{0,1\}^{n+1}}\)，使 \(g(x) \subset f_n(p(x))\)（对所有 \(x\)），且 \((g(x) \times g(x')) \cap \mathrm{A}_n\) 为空（对两个不同元素的每个偶 \((x,x')\)，属于 \(\{0,1\}^{n+1}\)）。然后定义映射 \(f_{n+1}\) 如下：对每个元素 \(x \in \{0,1\}^{n+1}\)，选取 \(g(x)\) 的一个非空开子集，其直径 \(\leqslant 2^{-n-1}\) 且其闭包包含于 \(g(x)\)。

对每个元素 \(x = (x_n)_{n \in \mathbf{N}}\)（属于 \(\{0,1\}^{\mathbf{N}}\)），集合的序列 \((f_n(x_0, \ldots, x_{n-1}))_{n \in \mathbf{N}}\) 是 P 的开子集的一个递减序列，其中每一个都包含后一个的闭包，且其直径趋于 0；因此这个集合序列的交退缩为一点（TG, II, p. 15），把它记作 \(f(x)\)。若两点 \(x\)、\(x'\)（属于 \(\{0,1\}^{\mathbf{N}}\)）满足 \(x_i = x'_i\)（对 \(i \leqslant n\)），则有 \(d(f(x), f(x')) \leqslant 2^{-n}\)。于是，映射 \(f: \{0,1\}^{\mathbf{N}} \to \mathbf{P}\) 连续。设 \(x\) 与 \(x'\) 是 \(\{0,1\}^{\mathbf{N}}\) 的两个不同元素。对 \(n \in \mathbf{N}\)（使 \((x_0, \ldots, x_n) \neq (x_0', \ldots, x_n')\)），开集 \(f_{n+1}(x_0, \ldots, x_n) \times f_{n+1}(x_0', \ldots, x_n')\) 与 \(A_n\) 不交（按 \(f_{n+1}\) 的定义），故偶 \((f(x), f(x'))\) 不属于 \(A_n\)。由此得出 \(f(x)\) 与 \(f(x')\) 对关系 R 不等价。于是，\(f\) 是单射的，且其像与 R 的每个等价类至多交于一点。命题由此得证。

